The design separates two tasks that matter for interpreting a correspondence system. The first is to retain at least one usable location, which we measure as candidate \emph{reach}. The second is to select that location when other candidates are present. A sparse set makes the first task hard. A larger set gives the scorer more opportunities and more competitors. To isolate those effects we vary the cap at fixed physical support, compare fit-only and diversified prefixes of equal length, and score the same centers with two different functions. This gives a paired account of gains from additional reach and changes in choices on tracks already reachable.

This separation also makes compute a design variable. An exhaustive lattice provides a useful accuracy reference, while the full diversified search limits how many local descriptors need to be computed. We evaluate both and measure local matching separately from shared scan initialization. The staging data change the institution and clinical setting while preserving an independently linked 3D reference. The screening tracks provide a replication with a different lesion population and acquisition setting.
